\begin{frame}[t]{\surf{surface 2 of 7}\ev{\tagO\,\tagP}Filesystem state shared by a snapshot}
\vspace{0pt}
{\small
\begin{tabularx}{\textwidth}{@{}L{2.8cm}L{4.1cm}Y@{}}
\toprule
State & Policy & Known mechanisms\\
\midrule
files, build outputs & pin by content digest & OBuilder takes a copy-on-write snapshot per build step (ZFS, btrfs)\\
workspace disk & copy-on-write layers & overlay layers (day10), DeltaBox's layered copy-on-write filesystem and chained overlaybd snapshots (DeepSeek DSec)\\
agent history & keep outside the sandbox & \\
randomness & reseed, or couple deliberately & \leavevmode{\color{Muted}see slide 38}\\
external services & mediate access & \leavevmode{\color{Muted}see slide 27}\\
\bottomrule
\end{tabularx}\par}
\claim{A snapshot captures only the state inside the sandbox boundary.}
\sources{\href{https://github.com/ocurrent/obuilder}{OBuilder}; \href{https://www.tunbury.org/2026/02/16/day10/}{day10}; \href{https://arxiv.org/abs/2605.22781}{DeltaBox}; \href{https://arxiv.org/abs/2606.19348}{DeepSeek-V4} \S5.2.5.}
\note{[0:45] Surface two is filesystem state. A snapshot, a saved copy of the sandbox, captures only the state inside that boundary. Each kind of state therefore needs its own policy. Build outputs are pinned by digest, and OBuilder snapshots every step. The workspace disk is copy-on-write, and DeepSeek's DSec already chains such snapshots. Randomness and external services come later. A filesystem snapshot cannot share memory, which the forking section addresses.}
\end{frame}
